\documentclass[11pt,a4paper]{article}
\usepackage[utf8]{inputenc}
\usepackage[T1]{fontenc}
\usepackage[margin=1in]{geometry}
\usepackage{amsmath}
\usepackage[
    colorlinks=true,  % Makes the links colored instead of boxed
    linkcolor=green,   % Color for internal links (e.g., table of contents)
    citecolor=blue,  % Color for citations
    urlcolor=cyan,    % Color for URLs/external links
]{hyperref}
\usepackage{amssymb}
\usepackage{graphicx}
\usepackage{amsthm}
\usepackage{amssymb}
\usepackage{physics,csquotes}

\newtheorem{theorem}{Theorem}
\newtheorem{lemma}{Lemma}

\title{Eigensystem of U(1) Covariant Lindbladians}
\author{Ashlin V Thomas}
\date{\today}

\begin{document}
\maketitle

\begin{abstract}
    In this work, we present an exact method for finding eigensystem of Lindbladians that are covariant under a $U(1)$ symmetry. By exploiting the block-diagonal structure induced by the symmetry, we derive recurrence relations for the coefficients of the left and right eigenoperators within each invariant subspace. 
    We provide explicit solutions for the cases where all jump operators have either positive or negative charges, constructing the corresponding eigenoperators using weighted directed graphs. 
    We map the problem of finding the eigenoperators to a combinatorial problem of counting weighted paths in these graphs, allowing for a systematic construction of the eigenoperators.
    The study reveals a complementary structure between the left and right eigenoperators, ensuring that their products yield Fock states. 
\end{abstract}

\section{$U(1)$ Covariant Lindbladians}

A Markovian quantum dissipative system is described via a master equation $\dot{\rho} = \mathcal{L} \rho$ for the density matrix $\rho$,
where the Liouvillian $\mathcal{L}$ is constrained to have the Lindblad or Gorini-Kossakowski-Sudarshan-Lindblad form. 
\begin{equation}
    \mathcal{L} \rho = -i[H, \rho] + \sum_{j} \gamma_j \left( L_j \rho L_j^\dagger - \frac{1}{2} \{ L_j^\dagger L_j, \rho \} \right) \label{Lindblad master equation}
\end{equation}
Here, $H$ is the system Hamiltonian, $L_j$ are the jump operators, and $\gamma_j$ are the associated decay rates. \\
In this work, we focus on Lindbladians that are covariant under a $U(1)$ symmetry. 
In particular, we consider generators that respect phase rotations generated by 
the number operator $\hat{N} = \hat{a}^\dagger \hat{a}$, 
characterized by the covariance conditions
\begin{equation}
    [\hat{H},\hat{N}] = 0, \qquad [\hat{N},L_j] = q_j\, L_j \label{U1 covariance}
\end{equation}
where $q_j \in \mathbb{Z}$ are the charges associated with each jump operator\cite{albert2014symmetries,gu2024spontaneous}. A large class of 
Since $\hat{H}$ commutes with $\hat{N}$ which has a non-degenerate spectrum, they share the common eigenbasis of Fock states:
\begin{equation}
    \hat{H} = \sum_{n=0}^{\infty} E_n \ket{n}\bra{n} \label{Hamiltonian spectrum}
\end{equation}
Further, we can obtain the spectral decomposition of the jump operators using eqn.(\ref{U1 covariance}):
\begin{align}
    N L_j   -  L_j N = q_j L_j \implies \sum_{m,n} (m-n - q_j) \bra{m} L_j \ket{n} \ket{m}\bra{n} = 0 
\end{align}
\text{Using the linear independence of the Fock basis, we get:}
\begin{align}
    \bra{m} L_j \ket{n} \neq 0 \quad \text{only if } m= n+ q_j \implies L_j = \begin{cases} \sum_{m=0}^{\infty} \alpha_{j,m} \ \ket{m}\bra{m-q_j}   &  \text{if } q_j < 0 \\ \sum_{m=0}^{\infty} \beta_{j,m} \ \ket{m+q_j}\bra{m} & \text{if } q_j > 0 \end{cases}                    \label{jump operator decomposition}
\end{align}
where $\alpha_{j,m} = \bra{m} L_j \ket{m-q_j}$ and $\beta_{j,m} = \bra{m+q_j} L_j \ket{m}$. \\

\section{Eigenspaces of the Linbladian}
Define an operator $\mathcal{N}: \mathcal{B}(\mathcal{H}) \to \mathcal{B}(\mathcal{H})$ acting on the space of bounded operators on the Hilbert space $\mathcal{H}$ as:
\begin{equation}
    \mathcal{N}(\hat{\rho}) = [\hat{N}, \hat{\rho}]
\end{equation}
\textit{\textbf{Lemma 2.1: }} Under the covariance conditions in eqn.(\ref{U1 covariance}), the Liouvillian $\mathcal{L}$ commutes with $\mathcal{N}$. \\
\textit{Proof: }    Using eqn.(\ref{Lindblad master equation}), we have:
\begin{align}
    \mathcal{N} (\mathcal{L} \rho) = & [\hat{N}, -i[H, \rho] + \sum_{j} \gamma_j ( L_j \rho L_j^\dagger - \frac{1}{2} \{ L_j^\dagger L_j, \rho \} ) ] \\
    = & -i [\hat{N}, H \rho - \rho H] + \sum_{j} \gamma_j \left( [\hat{N}, L_j \rho L_j^\dagger] - \frac{1}{2} [\hat{N}, L_j^\dagger L_j \rho] - \frac{1}{2} [\hat{N}, \rho L_j^\dagger L_j] \right)
\end{align}
In the first term, we can use the covariance condition $[\hat{H},\hat{N}] = 0$ and expanding the commutator in the following terms, we get 
$[\hat{N}, L_j \rho L_j^\dagger] = [\hat{N}, L_j] \rho L_j^\dagger + L_j [\hat{N}, \rho] L_j^\dagger + L_j \rho [\hat{N}, L_j^\dagger]$. Using the covariance conditions $[\hat{N}, L_j] = -q_j L_j$ and $[\hat{N}, L_j^\dagger] = q_j L_j^\dagger$, the first and last terms cancel out, leaving us with:
\begin{align}
    \mathcal{N} (\mathcal{L} \rho) = & -i [H, [\hat{N}, \rho]] + \sum_{j} \gamma_j \left( L_j [\hat{N}, \rho] L_j^\dagger - \frac{1}{2} \{ L_j^\dagger L_j, [\hat{N}, \rho] \} \right)  = \mathcal{L} (\mathcal{N}(\rho))
\end{align}
\textit{\textbf{Theorem 2.1: }} For any integer $l \in \mathbb{Z}$, the subspace $V_l = \text{span} \{ |n\rangle \langle n+l| : n \in \mathbb{N}_0 \}$ is invariant under the action of the Liouvillian $\mathcal{L}$. \\
\textit{Proof: } Consider an arbitrary operator $\rho \in V_l$, which can be expressed as:
\begin{equation}
    \rho = \sum_{n=0}^{\infty} \rho_n \ |n\rangle \langle n+l|
\end{equation}
Using the lemma, we have:
\begin{align}
    \mathcal{N}(\mathcal{L} \rho) = \mathcal{L}(\mathcal{N} \rho) = \mathcal{L}([\hat{N}, \rho]) = \mathcal{L} \left( \sum_{n=0}^{\infty} \rho_n (n - (n+l)) |n\rangle \langle n+l| \right) = -l \mathcal{L}(\rho)
\end{align}
This shows that $\mathcal{L} \rho$ is an eigenoperator of $\mathcal{N}$ with eigenvalue $-l$. Since the eigenspaces of $\mathcal{N}$ are precisely the subspaces $V_l$, it follows that $\mathcal{L} \rho \in V_l$. Thus, each subspace $V_l$ is invariant under the action of the Liouvillian $\mathcal{L}$. \hspace{4 mm} $\blacksquare$ \\
\textit{\textbf{Remark: }} A similar calculation shows that the adjoint Liouvillian $\mathcal{L}^\dagger$ also commutes with $\mathcal{N}$, and hence the subspaces $V_l$ are also invariant under the action of $\mathcal{L}^\dagger$. \\
Therefore, the left and right eigenoperators of $\mathcal{L}$ can be chosen to lie within these invariant subspaces $V_l$. This block-diagonal structure significantly simplifies the analysis of the Liouvillian's spectral properties, as we can study its action separately on each subspace $V_l$.

\section{Exact Diagonalization within Each Eigenspace}
\subsection{Right Eigenoperators}
Let us now focus on diagonalizing the Liouvillian $\mathcal{L}$ within each invariant subspace $V_l$. Let $R \in V_l$ be a right eigenoperator of the Liouvillian with eigenvalue $\lambda$ (i.e., $\mathcal{L} R = \lambda R$). Since $R \in V_l$, it can be expressed as:
\begin{equation}
    R = \sum_{n=0}^{\infty} r_n \ |n\rangle \langle n+l|
\end{equation}
Substituting this expression into the eigenvalue equation $\mathcal{L} R = \lambda R$ and using the explicit form of the Liouvillian from eqn.(\ref{Lindblad master equation}) alongwith the spectral decompositions in eqns.(\ref{Hamiltonian spectrum}) and (\ref{jump operator decomposition}), we obtain:
\begin{align}
    \lambda R = \mathcal{L} R = & -i [H, R] + \sum_{j} \gamma_j \left( L_j R L_j^\dagger - \frac{1}{2} L_j^\dagger L_j R - \frac{1}{2} R L_j^\dagger L_j \right) \\
    = & -i \sum_{n} r_n (E_n - E_{n+l}) |n\rangle \langle n+l| + \sum_{n,j} \gamma_j \alpha_{j,n - q_j} \alpha_{j,n+l - q_j}^* r_{n }\ |n - q_j \rangle \langle n + l - q_j| \nonumber \\
    & - \frac{1}{2} \sum_{n,j} \gamma_j r_n \left( |\alpha_{j,n-q_j}|^2 + |\alpha_{j,n+l-q_j}|^2 \right) |n\rangle \langle n+l| 
\end{align}
Taking inner products with the basis elements $\langle k |$ and $|k+l \rangle$, on both sides, we get the following recurrence relation for the coefficients $r_n$:
\begin{align}
    \lambda r_k = & -i r_k (E_k - E_{k+l}) + \sum_{j} \gamma_j \alpha_{j,k } \alpha_{j,k+l}^* r_{k + q_j} - \frac{1}{2} \sum_{j} \gamma_j r_k \left( |\alpha_{j,k-|q_j|}|^2 + |\alpha_{j,k+l-|q_j|}|^2 \right) \\
    \implies & \left( \lambda + i (E_k - E_{k+l}) + \frac{1}{2} \sum_{j} \gamma_j \left( |\alpha_{j,k-|q_j|}|^2 + |\alpha_{j,k+l-|q_j|}|^2 \right) \right) r_k = \sum_{j} \gamma_j \alpha_{j,k } \alpha_{j,k+l}^* r_{k + |q_j|} \label{right recurrence}
\end{align}
The above equation relates the coefficient $r_k$ to the coefficients $r_{k + q_j}$ for various $j$. To find the right eigenoperators explicitly, one needs to solve the above recurrence relation for the coefficients $r_k$. 

The above recurrence relation was derived specilaizing into the case of $l \geq 0$ and $sgn(q_j) = -1$. The recurrence relation can be stated in its full generality as follows:
\begin{equation}
    \left[ \lambda + i (E_m - E_{m+l}) + \frac{1}{2} \sum_{j} \gamma_j ( |\alpha_{j,m+q_j}|^2 + |\alpha_{j,m+l+q_j}|^2 ) \right] r_m = \sum_{j} \gamma_j \alpha_{j,m } \alpha_{j,m+l}^* r_{m - q_j} 
\end{equation}
where the sum on the right hand side is taken over all $j$ such that $m - q_j \geq 0$.


\subsection{Left Eigenoperators}
Before solving the recurrence relation for the right eigenoperators, let us setup a similar recurrence relation for the left eigenoperators. Let $L \in V_l$ be a left eigenoperator of the Liouvillian with eigenvalue $\lambda$ (i.e., $\mathcal{L}^\dagger L = \lambda^{*} L$). Since $L \in V_l$, it can be expressed as:
\begin{equation}
    L = \sum_{n=0}^{\infty} l_n \ |n\rangle \langle n+l|
\end{equation}
Substituting this expression into the eigenvalue equation $\mathcal{L}^\dagger L = \lambda^{*} L$ and using the explicit form of the adjoint Liouvillian, we obtain:
\begin{align}
    \lambda^{*} L = \mathcal{L}^\dagger L = & i [H, L] + \sum_{j} \gamma_j \left( L_j^\dagger L L_j - \frac{1}{2} L_j^\dagger L_j L - \frac{1}{2} L L_j^\dagger L_j \right) \\
    = & i \sum_{n} l_n (E_n - E_{n+l}) |n\rangle \langle n+l| + \sum_{n,j} \gamma_j \alpha_{j,n}^* \alpha_{j,n+l} l_n \ |n + q_j \rangle \langle n + l + q_j| \nonumber \\
    & - \frac{1}{2} \sum_{n,j} \gamma_j l_n \left( |\alpha_{j,n-q_j}|^2 + |\alpha_{j,n+l-q_j}|^2 \right) |n\rangle \langle n+l|
\end{align}
Taking inner products with the basis elements $\langle k |$ and $|k+l \rangle$, on both sides, we get the following recurrence relation for the coefficients $l_n$:
\begin{align}
    \lambda^{*} l_k = & i l_k (E_k - E_{k+l}) + \sum_{j} \gamma_j \alpha_{j,k-q_j}^* \alpha_{j,k+l-q_j} l_{k - q_j} - \frac{1}{2} \sum_{j} \gamma_j l_k \left( |\alpha_{j,k-q_j}|^2 + |\alpha_{j,k-l-q_j}|^2 \right) \\
    \implies & \left( \lambda^{*} - i (E_k - E_{k+l}) + \frac{1}{2} \sum_{j} \gamma_j \left( |\alpha_{j,k-q_j}|^2 + |\alpha_{j,k+l-q_j}|^2 \right) \right) l_k = \sum_{j} \gamma_j \alpha_{j,k - q_j}^* \alpha_{j,k+l - q_j} l_{k - q_j} \label{left recurrence}
\end{align}
Similar to the right eigenoperator case, the above equation relates the coefficient $l_k$ to the coefficients $l_{k - q_j}$ for various $j$ and one needs to solve this recurrence relation to find the left eigenoperators explicitly.\\
Solution of the recurrence relations in eqns.(\ref{right recurrence}) and (\ref{left recurrence}) largely depends on the charges $q_j$ of the jump operators. 
Solving these recurrence relations in full generality is challenging, but we can obtain explicit solutions in certain special cases. In the two subsequent sections, we present exact solutions for the cases where all jump operators have positive charge and where all jump operators have negative charge.



\section{Jump Operators with Positive Charge}
\textit{\textbf{Theorem 3.1: }} If all jump operators have positive charge (i.e., $q_j > 0$ for all $j$), then 
\begin{equation}
    \lambda_{m,l} = -i (E_m - E_{m+l}) - \frac{1}{2} \sum_{j} \gamma_j \left( |\alpha_{j,m+q_j}|^2 + |\alpha_{j,m+l+q_j}|^2 \right) \label{eigenvalue positive charge}
\end{equation}
is an eigenvalue of the Liouvillian $\mathcal{L}$ within the subspace $V_l$, for all $m \geq \max_{j} q_j $, assuming $\lambda_{m,l} \neq \lambda_{k,l}$ for $k \neq m$. \\
To construct the corresponding right and left eigenoperators, let us define a weighted directed graph $G_m$ whose vertices are $V = \mathbb{Z}_{\geq 0}$
and for each vertex $k$, there is a directed edge from vertex $k$ to vertex $k + q$ with some weight given by the following adjacency matrix:
\begin{gather}
    W_{k,k-q} = \frac{1}{\lambda_{m,l} - \lambda_{k,l}} \sum_{j, q_j=q} \gamma_j \bra{k}L_j\ket{k-q} \bra{k+l}L_j\ket{k+l-q}^*  \\
    W_{k,k+q} = \frac{1}{\lambda_{m,l}^* - \lambda_{k,l}^*} \sum_{j, q_j=q} \gamma_j \bra{k}L_j\ket{k+q}^* \bra{k+l}L_j\ket{k+l+q}
\end{gather}
Let $\mathcal{P}(i \to j)$ denote the set of all directed paths from vertex $i$ to vertex $j$ in the graph $G_m$. 
Define the weight of a path $W(P)$ as the product of the weights of the edges along the path. Then, the right eigenoperator $R_{m,l}$ and left eigenoperator $L_{m,l}$ corresponding to the eigenvalue $\lambda_{m,l}$ are given by:
\begin{equation}
    R_{m,l} = \sum_{k=0}^{m} \left( \sum_{P \in \mathcal{P}(k \to m)} W(P) \right) |k\rangle \langle k+l|, \quad L_{m,l} = \sum_{k=m}^{\infty} \left( \sum_{P \in \mathcal{P}(m \to k)} W(P) \right)^* |k\rangle \langle k+l|
\end{equation}
\textit{Proof: } To construct the right eigenoperator, we start from the recurrence relation in eqn.(\ref{right recurrence}). Setting $\lambda = \lambda_{m,l}$, we have:
\begin{align}
    \left( \lambda_{m,l} - \lambda_{k,l} \right) r_k = \sum_{j} \gamma_j \alpha_{j,k } \alpha_{j,k+l}^* r_{k + q_j}
\end{align}
We seek eigenvectors that satisfy $r_k = 0$ for $k > m$ and $r_m = 1$. Using these boundary conditions, we can recursively solve for $r_k$ for $k < m$:
\begin{itemize}
    \item For $k = m-1$, if there exists a jump operator with charge $q_j = 1$, or equivalently an edge from $m-1$ to $m$ in the graph $G_m$, we have:
    \begin{align}
        r_{m-1} = \frac{1}{\lambda_{m,l} - \lambda_{m-1,l}} \sum_{j, q_j=1} \gamma_j \alpha_{j,m-1} \alpha_{j,m-1+l}^* \ r_m = w_{m-1 \rightarrow m}  
    \end{align}
    If no such jump operator exists, then $r_{m-1} = 0$.
    \item For $k = m-2$, we have:
    \begin{align}
        r_{m-2} = & \frac{1}{\lambda_{m,l} - \lambda_{m-2,l}} \sum_{j, q_j = 2} \gamma_j \alpha_{j,m-2} \alpha_{j,m-2+l}^* \ r_m + \frac{1}{\lambda_{m,l} - \lambda_{m-2,l}} \sum_{j, q_j = 1} \gamma_j \alpha_{j,m-2} \alpha_{j,m-2+l}^* \ r_{m-1} 
    \end{align}  
    \begin{gather}
        =   w_{m-2 \rightarrow m} + \frac{1}{\lambda_{m,l} - \lambda_{m-2,l}} \sum_{j, q_j = 1} \gamma_j \alpha_{j,m-2} \alpha_{j,m-2+l}^* \ w_{m-1 \rightarrow m} \\  = w_{m-2 \rightarrow m} + w_{m-2 \rightarrow m-1}  w_{m-1 \rightarrow m} 
        =  \sum_{P \in \mathcal{P}(m-2 \to m)} W(P)
    \end{gather}
    \item Continuing this process recursively, we find that for any $k < m$:
    \begin{align}
        r_k = \sum_{P \in \mathcal{P}(k \to m)} W(P)
    \end{align}
\end{itemize}
By construction, we have $r_m = 1$ and hence the the operator is non-zero. Therefore, the right eigenoperator is given by:
\begin{equation}
    R_{m,l} = \sum_{k=0}^{m} \left( \sum_{P \in \mathcal{P}(k \to m)} W(P) \right) |k\rangle \langle k+l|
\end{equation}
To construct the left eigenoperator, we start from the recurrence relation in eqn.(\ref{left recurrence}). Setting $\lambda = \lambda_{m,l}$, we have:
\begin{align}
    \left( \lambda_{m,l}^* - \lambda_{k,l}^* \right) l_k = \sum_{j} \gamma_j \alpha_{j,k - q_j}^* \alpha_{j,k+l - q_j} l_{k - q_j}
\end{align}
We seek eigenvectors that satisfy $l_k = 0$ for $k < m$ and $l_m = 1$. Using these boundary conditions, we can recursively solve for $l_k$ for $k > m$:
\begin{itemize}
    \item For $k = m+1$, if there exists a jump operator with charge $q_j = 1$, or equivalently an edge from $m$ to $m+1$ in the graph $G_m$, we have:
    \begin{align}
    l_{m+1} = \frac{1}{ \left( \lambda_{m,l}^* - \lambda_{m+1,l}^* \right)} \sum_{j, q_j=1} \gamma_j \alpha_{j,m}^* \alpha_{j,m+l}  \implies l_{m+1} = w_{m \rightarrow m+1}^*
    \end{align}
    If no such jump operator exists, then $l_{m+1} = 0$.
    \item For $k = m+2$, we have:
    \begin{gather}
        \left( \lambda_{m,l}^* - \lambda_{m+2,l}^* \right) l_{m+2} =  \sum_{j, q_j = 2} \gamma_j \alpha_{j,m}^* \alpha_{j,m+l}  + \sum_{j, q_j = 1} \gamma_j \alpha_{j,m+1 }^* \alpha_{j,m+1+l } l_{m+1} \\
        \implies l_{m+2} =  w_{m \rightarrow m+2}^* + \frac{1}{ \left( \lambda_{m,l}^* - \lambda_{m+2,l}^* \right)} \sum_{j, q_j = 1} \gamma_j \alpha_{j,m+1}^* \alpha_{j,m+1+l } w_{m \rightarrow m+1}^* \\
        =  w_{m \rightarrow m+2} + w_{m \rightarrow m+1}^* w_{m+1 \rightarrow m+2} = \sum_{P \in \mathcal{P}(m \to m+2)} W(P)^* 
    \end{gather}
    \item Continuing this process recursively, we find that for any $k > m$:
    \begin{align}
        l_k = \sum_{P \in \mathcal{P}(m \to k)} W(P)^*
    \end{align}
\end{itemize}
By construction, we have $l_m = 1$ and hence the the operator is non-zero. Therefore, the left eigenoperator is given by:
\begin{equation}
    L_{m,l} = \sum_{k=m}^{\infty} \left( \sum_{P \in \mathcal{P}(m \to k)} W(P) \right)^* |k\rangle \langle k+l|
\end{equation}
This completes the proof. \hspace{4 mm} $\blacksquare$ \\

\section{Jump Operators with Negative Charge}
\textit{\textbf{Theorem 4.1: }} If all jump operators have negative charge (i.e., $q_j < 0$ for all $j$), then $\lambda_{m,l}$ as defined in eqn.(\ref{eigenvalue positive charge}) is an eigenvalue of the Liouvillian $\mathcal{L}$ within the subspace $V_l$, for all 
$m \geq \max_{j} |q_j| $, assuming $\lambda_{m,l} \neq \lambda_{k,l}$ for $k \neq m$. \\
The left and right eigenoperators corresponding to the eigenvalue $\lambda_{m,l}$ can be constructed similarly as in Theorem 3.1, but with the roles of the left and right eigenoperators interchanged. Specifically, 
\begin{align}
    R_{m,l} = \sum_{k=m}^{\infty} \left( \sum_{P \in \mathcal{P}(m \to k)} W(P) \right)^* |k\rangle \langle k+l|, \quad L_{m,l} = \sum_{k=0}^{m} \left( \sum_{P \in \mathcal{P}(k \to m)} W(P) \right) |k\rangle \langle k+l|
\end{align}
\textit{Proof: } The proof follows similarly to that of Theorem 3.1, but with the roles of the left and right eigenoperators interchanged. The recurrence relations in eqns.(\ref{right recurrence}) and (\ref{left recurrence}) now relate coefficients in the opposite direction due to the negative charges of the jump operators. By imposing appropriate boundary conditions and recursively solving the relations, we arrive at the stated forms for the right and left eigenoperators. \hspace{4 mm} $\blacksquare$ \\

\section{Structure of $L^{\dagger}R$ and $L R^{\dagger}$}
In the two previous sections, we constructed the left and right eigenoperators corresponding to each eigenvalue $\lambda_{m,l}$ for the cases where all jump operators have positive or negative charge. 
In both cases, we found that if the right eigenoperator has zero coefficients beyond a certain index $m$, then the left eigenoperator has non-zero coefficients only beyond that index $m$, and vice-versa.
This complementary structure ensures that the products $L^{\dagger} R$ and $L R^{\dagger}$ are always fock states.

\section{Examples}
Now, we illustrate the above constructions with some explicit examples.
\subsection{Single mode bosonic oscillator with single photon loss}
Consider a single mode bosonic oscillator with Hamiltonian $H = \omega a^{\dagger} a$ and a single jump operator $L = a$ corresponding to photon loss. 
Note that $[a , N] = a$ and hence the jump operator has charge $q = 1$. Using eqn.(\ref{jump operator decomposition}), we have $\alpha_{n} = \langle n | \hat{a} | n+1 \rangle = \sqrt{n+1}$.
Since $H = N$, $E_n = \omega n$. Using eqn.(\ref{eigenvalue positive charge}), the eigenvalues of the Liouvillian within the subspace $V_l$ are given by:
\begin{equation}
    \lambda_{m,l} = +i \omega l - \frac{\gamma}{2} (2 m + l)
\end{equation}
where $m \geq 0$. The corresponding right and left eigenoperators can be constructed using Theorem 3.1. 
Since we have only one jump operator with charge $q=1$, the graph $G_m$ has edges from vertex $k$ to vertex $k+1$ with weights:
\begin{equation}
    w_{k \to k+1} = \frac{\gamma \sqrt{(k+1)(k+l+1)}}{\lambda_{m,l} - \lambda_{k,l}} = \frac{ \sqrt{(k+1)(k+l+1)}}{  (k - m)}
\end{equation}
The only directed path from vertex $k$ to vertex $m$ is the path that goes through all intermediate vertices, and its weight is given by:
\begin{equation}
    W(P) = \prod_{s=k}^{m-1} w_{s \to s+1} = \prod_{s=k}^{m-1} \frac{ \sqrt{(s+1)(s+l+1)}}{  (s - m)} = (-1)^{m-k} \sqrt{\frac{m! (m+l)!}{k! (k+l)!}} \frac{1}{(m-k)!}
\end{equation}
Therefore, the right eigenoperator is given by:
\begin{equation}
    R_{m,l} = \sum_{k=0}^{m} W(P) |k\rangle \langle k+l| = \sum_{k=0}^{m} (-1)^{m-k} \sqrt{\frac{m! (m+l)!}{k! (k+l)!}} \frac{1}{(m-k)!} |k\rangle \langle k+l|
\end{equation}
Similarly, the only directed path from vertex $m$ to vertex $k$ is the path that goes through all intermediate vertices, and its weight is given by:  
\begin{equation}
    W(P) = \prod_{s=m}^{k-1} w_{s \to s+1} = \prod_{s=m}^{k-1} \frac{ \sqrt{(s+1)(s+l+1)}}{  (s - m)} = \sqrt{\frac{k! (k+l)!}{m! (m+l)!}} \frac{1}{(k-m)!}
\end{equation}
Therefore, the left eigenoperator is given by:
\begin{equation}
    L_{m,l} = \sum_{k=m}^{\infty} W(P)^* |k\rangle \langle k+l| = \sum_{k=m}^{\infty} \sqrt{\frac{k! (k+l)!}{m! (m+l)!}} \frac{1}{(k-m)!} |k\rangle \langle k+l|
\end{equation}
The constructed left and right eigenoperators agree with the known results for the eigensystem of the single photon loss Lindbladian.



\bibliography{refs.bib}
\bibliographystyle{plain}
\end{document}
